\documentclass[a4paper,10pt]{report}

\def\packagepath{../../preambule}   % path du package princial
\usepackage{\packagepath/preambule}    % utilisation du fichier de configuration

\def\level{Première }              % Classe
\def\course{NSI}              % Matière
\def\eval{Codage des entiers relatifs}

\def\visibleornot{visible}    % visible or invisible
\def\documentpath{./Sources_Latex}
\def\date{Mardi 08 Avril 2025}
\renewcommand{\arraystretch}{1}  % Ecart dans les tableau

\def\ptsexoA{3}
\def\ptsexoB{3}
\def\ptsexoC{2}
\def\ptsexoD{2}
\def\ptstotal{\ptsexoA+\ptsexoB}

\usepackage{hyperref}

\usetikzlibrary{shapes, arrows, positioning}
\usetikzlibrary{shapes.geometric, arrows, positioning, decorations.pathreplacing}

\tikzstyle{etat} = [draw, rounded corners=15pt, minimum width=2.5cm, minimum height=1.2cm, text centered, font=\bfseries]
\tikzstyle{nouveau} = [etat, fill=blue!40]
\tikzstyle{pret} = [etat, fill=green!60]
\tikzstyle{elu} = [etat, fill=yellow!60]
\tikzstyle{bloque} = [etat, fill=red!60]
\tikzstyle{termine} = [etat, fill=white]
\tikzstyle{fleche} = [thick, ->, >=stealth]

\usepackage{float} % Permet d'utiliser [H]
\usepackage[T1]{fontenc}
\begin{document}

%%%%%%%%%%%%%%%%%%%%%%%%%%%%%%%%%%%%%%%%%%%%%%%%%%ù

%\PageGardeSujetBac[Matiere = Numérique et Sciences Informatiques,
%                  Session = 2025,
%                  AffJour=false,
%                  Duree = {3 heures 30},
%                  DernierePage = \pageref{LastPage},
%                  ModeExamen=false
%                  ]
\renewcommand{\labelitemi}{\textbullet} %pour éviter les tirets dans les "itemize" qui apportent confusion avec le signe moins.
%% Début page de garde style BAC

   %%%%%%%%%%%%%%%%%%%%%%%%%%%%%%%%%%%%%%%%%%%%%%%%%%%%%%%%
%\PageGardeSujetBac[clés]

\pagestyle{DS_FP}          % Feuille de style fancy
\NomPrenomNote{}


\exods{\ptsexoA}

Répondre aux questions suivantes. Aucune justification n'est demandée.

\begin{enumerate}
    \item Avec 8 bits, combien de valeurs différentes peut-on coder? \dotfill
    \item Quelles sont les valeurs entières naturelles minimale et maximale codable sur 8 bits? \dotfill
    \item Quelles sont les valeurs entières relatives minimale et maximale codable sur 8 bits? \dotfill
\end{enumerate}


\exods{\ptsexoB}

\begin{enumerate}
    \item Ecrire l'algorithme de conversion d'un entier naturel compris entre 0 et 255 en binaire sur 8 bits. Vous pouvez écrire cet algorithme en pseudo-code ou en Python.
    
    \begin{answer}{1}{4}\end{answer}

    \item Donner l'algorithme de conversion d'un entier relatif en binaire sur 8 bits.
    
    Vous citerez les différentes étapes de l'algorithme sans le traduire en Python.
    \begin{answer}{1}{4}\end{answer}

    \item Par la méthode que vous avez explicité précédemment (question 2.), convertir les valeurs suivantes en binaire sur 8 bits:
    
    \begin{answer}{1}{7}
    
        \begin{minipage}[t]{0.33\linewidth}
            \centering{$-24$}
        \end{minipage}
        \begin{minipage}[t]{0.33\linewidth}
            \centering{$-124$}
        \end{minipage}
        \begin{minipage}[t]{0.33\linewidth}
            \centering{$9$}
        \end{minipage}
    
    \end{answer}

\end{enumerate}

\newpage

\exods{\ptsexoC}

On suppose que l'on dispose des fonctions suivantes en python:



\begin{CodePythontexAlt}[Largeur=1\linewidth,Centre,PremLigne=1]{}
def conversion_entier_binaire(val_entiere:int)->str:
    ''' Convertit un entier naturel en binaire sur 8 bits. 
    val_entiere : int : entier naturel compris entre 0 et 255
    return : str : représentation binaire sur 8 bits de l'entier naturel'''
    # Code non demandé
    return val_binaire    

assert conversion_entier_binaire(34) == '00100010'

def inverse_bit(bit:str)->str:
    ''' Inverse un bit 
    bit : str : bit à inverser
    return : str : bit inversé'''
    # Code non demandé
    return bit_inverse

assert inverse_bit('0') == '1'
assert inverse_bit('1') == '0'

def additionne_binaire(binaire1:str, binaire2:str)->str:
    ''' Additionne deux entiers binaires sur 8 bits
    binaire1 : str : premier entier binaire
    binaire2 : str : second entier binaire
    return : str : somme des deux entiers binaires '''
    # Code non demandé
    return somme_binaire

assert additionne_binaire('00000001','00000001') == '00000010'
assert additionne_binaire('11111111','00000001') == '00000000'

\end{CodePythontexAlt}

Ecrire la fonction \texttt{conversion\_relatif\_binaire} qui convertit un entier relatif en binaire sur 8 bits. Vous utiliserez les fonctons précédentes pour créer votre fonction. 
    
    \begin{answer}{1}{8.5}
    \begin{verbatim}
def conversion_relatif_binaire(val_entiere:int)->str:
    ''' Convertit un entier relatif en binaire sur 8 bits. 
    val_entiere : int : entier relatif compris entre -128 et 127
    return : str : représentation binaire sur 8 bits de l'entier relatif'''
    \end{verbatim}
    \end{answer}


\exods{\ptsexoD}


Entourer la ou les bonne(s) réponse(s)

\begin{enumerate}
    \item Que peut-on dire de l'entier relatif $10100011$ codé sur 8 bits en complément à deux ?
    \begin{enumerate}[A)]
        \begin{minipage}[t]{0.24\linewidth}
        \item Cet entier est pair 
        \end{minipage}\hfill
        \begin{minipage}[t]{0.24\linewidth}
            \item Cet entier est négatif
        \end{minipage}\hfill
        \begin{minipage}[t]{0.24\linewidth}
            \item Cet entier est positif
        \end{minipage}\hfill
        \begin{minipage}[t]{0.24\linewidth}
            \item Cet entier vaut $-12_{10}$
        \end{minipage}
    \end{enumerate}

    \medskip

    \item Si $1011$ est un entier signé sur 4 bits en complément à 2, alors sa valeur décimale est:
        \begin{enumerate}[A)]
        \begin{minipage}[t]{0.24\linewidth}
        \item $-7$
        \end{minipage}\hfill
        \begin{minipage}[t]{0.24\linewidth}
        \item $-5$
        \end{minipage}\hfill
        \begin{minipage}[t]{0.24\linewidth}
        \item $5$       
        \end{minipage}\hfill
        \begin{minipage}[t]{0.24\linewidth}
        \item $7$
        \end{minipage}
    \end{enumerate}
\end{enumerate}



\end{document}








